% Last updated: 2025/04/23

\subsection{Physician Training Pipeline}

\subsubsection{Medical School}
A physician's medical education begins in medical school.\footnote{In the US, there are both allopathic medical schools, which grant MDs, and osteopathic medical schools, which grant DOs. In this paper, I focus on MDs allopathic medical schools since only 11.4\% of US medical school graduates are DOs, with an even smaller percent of surgeons being DOs (\cite{noauthor_us_nodate}). Any mention of medical school henceforth will refer to allopathic medical schools.} During the four years of medical school, students first study a core curriculum through classes, then attain clinical skills through clinical rotations. Core rotations include family medicine, internal medicine, neurology, obstetrics and gynecology (OBGYN), pediatrics, psychiatry, and surgery. These take place at teaching hospitals affiliated with the medical school. For example, Emory University School of Medicine students rotate at Emory University Hospital as well as the nearby Grady Memorial Hospital, among others. During each rotation, medical students spend 4-8 weeks as a member of a medical team and improve on their clinical skills in a patient-facing setting. Rotations are designed to allow students to gain practical experience in evaluating and diagnosing real patients and performing simple procedures. They also play a large role in a medical student's specialty choice. Rotations are often the first time medical students have contact with patients. Observing how senior physicians diagnose and treat patients is foundational to how these future physicians develop habits in patient care. 

\subsubsection{Residency and Fellowship}
Additional physician training that takes place after medical school is collectively known as graduate medical education (GME) and includes both residency and fellowship. Surgeons typically spend 5-7 years in graduate medical training, depending on their surgical subspecialty. For example, becoming a general surgeon requires a 5-year general surgery residency, while becoming a vascular surgeon requires a 2-year vascular surgery fellowship in addition to a 5-year general surgery residency. This next stage of medical training is a hands-on learning experience that all physicians undergo before beginning as an independent practitioner. The first year of residency is spent as an intern who can only practice medicine under the supervision of an attending physician. This allows newly graduated doctors to learn on the job from practicing physicians in their specialty. Residents beyond intern year and fellows have more responsibility over their patients and further develop their own practice style through learned practice experience. 

\subsection{Physician Migration}

\subsubsection{Moving to Residency: The Match}
Individuals are typically matched to residency and fellowship programs by the National Resident Matching Program (NRMP), also known as the Match.\footnote{Some specialties use other match programs. These include ophthalmology (San Francisco (SF) Match), plastic surgery (SF Match), and urology (American Urological Association Match).} During the Match process, individuals first file their applications. The programs review the applications and select applicants to interview. Interviewed applicants and programs then create rank order lists, which list preference rankings of programs and applicants, respectively. A matching algorithm allocates applicants to programs based on the rank order lists and the number of available spots. Due to this matching process, individuals may not place at their top-choice program and thus do not have full control of where they are located for training.

\subsubsection{Moving to Independent Practice}
After their final year of GME, physicians start practicing independently in their specialty. A physician's final job selection depends on a variety of factors, including salary, practice type and size, and location. In particular, there is a documented location retention effect after GME.\footnote{There does not seem to be a similar retention effect for medical school. Only 15\% of my sample both attend medical school and currently practice in the same market.} According to the Association of American Medical Colleges (AAMC), 58.6\% of physicians who completed residency between 2014 and 2023 continued to practice in the state they completed residency in 2024 (\cite{noauthor_table_nodate}). Licensing likely plays a role in this home effect, as a physician usually becomes fully-licensed in the state of their residency program after intern year. The cost of becoming licensed in another state may result in doctors staying in the same area. The Interstate Medical Licensure Compact is a recent initiative that makes it easier for physicians to gain licensure in other states, mitigating the residency retention effect. 

\subsection{Physician Practice Intensity}
When physicians initially practice independently, their treatment decisions are influenced by both demand-side factors, e.g., patient demand and insurance, and supply-side factors, e.g., financial incentives and institutional norms, as well as their own treatment intensity. Treatment intensity can be generalized as a physician's beliefs about treatment. In the context of end-of-life care, \textcite{cutler2019} introduce two distinct practice styles: ``cowboys," who consistently recommend intensive treatment, and ``comforters," who consistently recommend palliative care. Most other papers defined practice intensity in a continuous manner, computing it as the propensity to perform a certain operation, such as c-sections. These differences in practice intensity has been shown to directly impact healthcare expenditures. \textcite{cutler2019} and \textcite{badinski2023} find that about a third of geographic variation in end-of-life spending and overall healthcare spending, respectively, is due to differences in physician practice style. The regional development of different practice styles is a key contributing factor to continued variation in healthcare expenditures. 

The literature presents strong evidence demonstrating that practice intensity continues to develop after a physician begins independent practice. Treatment patterns adjust to both changes in reimbursement rates and changes in new medical research (\cite{clemens2014}; \cite{decicca_how_2024}), and they are also responsive to a physician's peers in their practice environment. \textcite{epstein2009} find that OBGYNs tend to increase c-section rates in response to increases in rates among peers at the same hospital, and \textcite{molitor2018} finds that cardiologists who move from one market to another adopts the practice intensity of their destination market immediately after a move, changing by 0.63\% for every 1\% change in market intensity. These changes demonstrate how physicians learn and update their beliefs as they continue working. This updating is relative to a baseline that is established during a physician's time in training.

The influence of the factors that shape that baseline practice intensity are less clear. Before practicing independently, physicians undergo a lengthy training period consisting of multiple stages. Common sense dictates that this formal training period has a large effect on how a physician practices after training, but the literature has yet to quantify this effect in a causal way. For example, \textcite{epstein2009} find that only 2\% of variation in c-section rates can be attributed to residency program through the use of residency program fixed effects. Other papers in this area focus primarily on how the ranking of medical school and residency programs impact treatment intensity and healthcare utilization. \textcite{schnell_addressing_2018} find that physicians who attended lower-ranked medical schools tend to prescribe more opioids compared to those who attended higher-ranked medical schools. \textcite{doyle_jr_returns_2010} find that patients assigned to care teams from high-ranking residency programs incur lower costs. These papers show that training has an effect on how physicians practice, but the pathway is unclear. 

While rank is a good proxy for quality of education, it does not say anything about the style of care that new physicians are being taught. It's possible that education from higher-ranked institutions focus on new techniques and procedures, resulting in graduating classes who tend to treat patients more intensely. It's also possible that higher-ranked institutions choose good candidates for surgery more conservatively, resulting in graduating classes who tend to treat patients less intensely. In order to study how training impacts a physician's intensity more directly, I use a new measure approximating what medical students learn to be the ``standard of care" by averaging physician intensity at the market level to capture what is being taught in medical school. This allows me to test whether a surgeon's experience in medical school clinical rotations plays a large role in how they currently practice. In other words, I am able to test if graduating classes of medical schools practice similarly to their alma mater. 

